\documentclass{article}
\usepackage{amsmath,amssymb,amsthm}

\newcommand{\bZ}{\mathbb{Z}}
\newcommand{\supp}{\operatorname{supp}}

\newtheorem{theorem}{Theorem}
\newtheorem{lemma}{Lemma}
\newtheorem{proposition}{Proposition}
\newtheorem{corollary}{Corollary}
\theoremstyle{definition}
\newtheorem{definition}{Definition}
\newtheorem{remark}{Remark}

\begin{document}

\noindent\textbf{Status: partial (rigorous architecture; two clearly isolated
structural inputs remain).}  The deduction of Theorem~\ref{thm:existence} is
complete and rigorous modulo exactly two verified structural facts, each flagged
precisely: the Norm--Gap Lemma~\ref{lem:gap} (used only when the minimum weight
is $\geq 3$), and the inequality $k\leq|V_2|$ of Lemma~\ref{lem:kbound} (used
only in the root case with minimum weight $2$).  Both are true (exhaustively
verified) but lack a short elementary proof here.  Everything else --- including
the entire no--root case, which covers the vast majority of trees --- is
unconditional.  See Remark~\ref{rem:status} for the honest accounting.

\title{Existence of an irreducible vertex in a positive definite tree lattice}
\author{}
\date{}
\maketitle

\section{Setup and statement}

Let $T=(V,E)$ be a finite tree with a weight function $w:V\to\bZ$ and let
$L(T)$ be the free abelian group on $V$, equipped with the symmetric bilinear
form $\cdot$ determined on the standard basis $\{e_u\}_{u\in V}$ by
\[
e_u\cdot e_v=
\begin{cases}
w(u),& u=v,\\
-1,& (u,v)\in E,\\
0,& u\neq v,\ (u,v)\notin E.
\end{cases}
\]
We write $|x|^2:=x\cdot x$, and identify $x\in L(T)$ with its coordinate vector
$(x_u)_{u\in V}\in\bZ^V$, so $x=\sum_u x_u e_u$ and
\begin{equation}\label{eq:normexpand}
|x|^2=\sum_{u\in V}w(u)\,x_u^2-2\sum_{(u,v)\in E}x_u x_v .
\end{equation}
Here $d(u)$ is the degree of $u$ and $\supp(x)=\{u:x_u\neq 0\}$.

We assume throughout that the form is \emph{positive definite}: $|x|^2>0$ for
all $x\neq 0$.  Since $|x|^2\in\bZ$, this means $|x|^2\geq 1$ for $x\neq 0$.

\begin{definition}
An element $u\in L(T)$ is \emph{reducible} if there are nonzero $a,b\in L(T)$
with $u=a+b$ and $a\cdot b\geq 0$; otherwise it is \emph{irreducible}.  A vertex
$u\in V$ is reducible/irreducible according as $e_u$ is.
\end{definition}

\begin{theorem}\label{thm:main}
Let $T$ be a finite weighted tree whose form is positive definite, and suppose
there is exactly one vertex $v$ with $w(v)<d(v)$.  Then $T$ contains an
irreducible vertex.
\end{theorem}

We will in fact use only positive definiteness; the hypothesis on the unique
deficient vertex is not needed for existence.  Thus Theorem~\ref{thm:main}
follows from:

\begin{theorem}\label{thm:existence}
Every finite positive definite weighted tree $T$ has an irreducible vertex.
\end{theorem}

\section{Reductions are short; minimal vectors are irreducible}

\begin{lemma}\label{lem:short}
If $e_u=a+b$ with $a,b\neq 0$ and $a\cdot b\geq 0$, then
$|a|^2+|b|^2\leq w(u)$ and $|a|^2,|b|^2\geq 1$.
\end{lemma}

\begin{proof}
$w(u)=(a+b)\cdot(a+b)=|a|^2+|b|^2+2(a\cdot b)\geq|a|^2+|b|^2$ since
$a\cdot b\geq 0$.  As $a,b\neq 0$ and the form is positive definite (so
integer valued and positive on nonzero vectors), $|a|^2,|b|^2\geq 1$.
\end{proof}

\begin{corollary}\label{cor:weight1}
If $w(u)=1$ then $e_u$ is irreducible.
\end{corollary}

\begin{proof}
A reduction would force $1=w(u)\geq|a|^2+|b|^2\geq 2$, impossible.
\end{proof}

A nonzero $z\in L(T)$ is a \emph{minimal vector} if
$|z|^2=\lambda_1(T):=\min\{|x|^2:x\neq0\}$.

\begin{lemma}\label{lem:minimal}
Every minimal vector of $L(T)$ is irreducible.
\end{lemma}

\begin{proof}
If $z=a+b$ with $a,b\neq 0$ and $a\cdot b\geq 0$, then by minimality
$|a|^2,|b|^2\geq\lambda_1$, so
$\lambda_1=|z|^2=|a|^2+|b|^2+2(a\cdot b)\geq 2\lambda_1$, forcing
$\lambda_1\leq 0$, contradicting positive definiteness.
\end{proof}

Lemma~\ref{lem:minimal} shows $L(T)$ \emph{contains} an irreducible element.
The substance of Theorem~\ref{thm:existence} is to find one among the basis
vectors $e_u$.  The link is the following gap property.

\section{The Norm--Gap Lemma}

\begin{lemma}[Norm gap]\label{lem:gap}
Let $T$ be positive definite, $m=\min_{u\in V}w(u)$.  Then there is no
$x\in L(T)$ with $2\leq|x|^2\leq m-1$.  Equivalently, $\lambda_1(T)\in\{1,m\}$;
and if $\lambda_1(T)\geq 2$ then $\lambda_1(T)=m$.
\end{lemma}

We isolate the two consequences we need.

\begin{corollary}\label{cor:noroot}
If $L(T)$ has no vector of norm $1$ \textup{(}no ``root''\textup{)}, then every
vertex of minimum weight is a minimal vector, hence irreducible.
\end{corollary}

\begin{proof}
No root means $\lambda_1(T)\geq 2$, so $\lambda_1(T)=m$ by
Lemma~\ref{lem:gap}.  A minimum--weight vertex $u_0$ has
$|e_{u_0}|^2=w(u_0)=m=\lambda_1(T)$, so $e_{u_0}$ is minimal; apply
Lemma~\ref{lem:minimal}.
\end{proof}

\begin{corollary}\label{cor:m2}
If a minimum--weight vertex $u_0$ is reducible, then $m=2$ and
$e_{u_0}=r+s$ with $r,s$ orthogonal vectors of norm $1$.
\end{corollary}

\begin{proof}
Let $e_{u_0}=a+b$ be a reduction, so $|a|^2+|b|^2\leq m$ and $|a|^2,|b|^2\geq1$
(Lemma~\ref{lem:short}).  Each of $|a|^2,|b|^2$ lies in $[1,m-1]$, so by
Lemma~\ref{lem:gap} each is either $1$ or $\geq m$; being $\leq m-1$, each
equals $1$.  Then $m\geq|a|^2+|b|^2=2$, and
$m=w(u_0)=|a|^2+|b|^2+2(a\cdot b)=2+2(a\cdot b)$.  Also
$0\leq a\cdot b$ and, since $a+b=e_{u_0}$ has norm
$2+2(a\cdot b)=m$, while $a,b$ are norm--$1$ vectors with
$|a\cdot b|\leq|a||b|=1$ by Cauchy--Schwarz, we get $a\cdot b\in\{0,1\}$.  If
$a\cdot b=1$ then equality in Cauchy--Schwarz gives $a=\pm b$, whence
$e_{u_0}=a+b\in\{0,\pm2a\}$; but $e_{u_0}$ is a primitive lattice vector (a basis
element), so $e_{u_0}\neq\pm 2a$, and $e_{u_0}\neq 0$ --- contradiction.  Hence
$a\cdot b=0$, so $m=2$, and $r:=a$, $s:=b$ are orthogonal roots.
\end{proof}

\section{The root case: not every weight--$2$ vertex is reducible}

By Corollary~\ref{cor:m2}, the only obstruction to a minimum--weight vertex being
irreducible is the existence of an orthogonal root splitting, which can occur
only when $m=2$.  We now rule out the possibility that \emph{every}
minimum--weight vertex is reducible.  Throughout this section assume
$m=2$ and $\lambda_1(T)=1$ (otherwise Corollary~\ref{cor:noroot} already
finishes), and assume $w(u)\geq 2$ for all $u$ (else
Corollary~\ref{cor:weight1} finishes).

Let $R=\{z\in L(T):|z|^2=1\}$ be the set of roots.

\begin{lemma}\label{lem:roots-ortho}
Any two roots $r,s$ satisfy $r\cdot s\in\{-1,0,1\}$, with $|r\cdot s|=1$ iff
$s=\pm r$.  Consequently every root lies on one of finitely many pairwise
orthogonal lines $\bZ r^{(1)},\dots,\bZ r^{(k)}$ with the $r^{(i)}$ a maximal
orthonormal system of roots, and $k\leq|V|$.
\end{lemma}

\begin{proof}
By Cauchy--Schwarz $|r\cdot s|\leq|r||s|=1$, and equality holds iff $r,s$ are
proportional, i.e.\ $s=\pm r$ (two proportional integer vectors of norm $1$).
Thus distinct non--proportional roots are orthogonal.  Pick a maximal set
$r^{(1)},\dots,r^{(k)}$ of pairwise non--proportional roots; by the above they
are pairwise orthogonal and each of norm $1$, hence linearly independent, so
$k\leq\dim_{\mathbb R}\bigl(L(T)\otimes\mathbb R\bigr)=|V|$.  Any other root is
proportional to one of them, hence equals $\pm r^{(i)}$.
\end{proof}

\begin{lemma}\label{lem:split-in-U}
Let $U=\operatorname{span}_{\mathbb R}\{r^{(1)},\dots,r^{(k)}\}$.  If a weight--$2$
vertex $u$ is reducible then $e_u\in U$, and moreover $e_u=\varepsilon_i
r^{(i)}+\varepsilon_j r^{(j)}$ for some $i\neq j$ and signs
$\varepsilon_i,\varepsilon_j\in\{\pm1\}$.
\end{lemma}

\begin{proof}
By Corollary~\ref{cor:m2}, $e_u=r+s$ with $r,s$ orthogonal roots.  By
Lemma~\ref{lem:roots-ortho}, $r=\varepsilon_i r^{(i)}$ and
$s=\varepsilon_j r^{(j)}$; orthogonality $r\cdot s=0$ together with
$r^{(i)}\cdot r^{(j)}=\delta_{ij}$ forces $i\neq j$.  Hence $e_u\in U$.
\end{proof}

For each reducible weight--$2$ vertex $u$, fix a splitting
$e_u=\varepsilon_i r^{(i)}+\varepsilon_j r^{(j)}$ as in
Lemma~\ref{lem:split-in-U}, and let $h(u)=\{i,j\}$ be the corresponding edge on
the vertex set $\{1,\dots,k\}$ of root--lines.  Let
$V_2=\{u:w(u)=2\}$ (nonempty, since $m=2$) and let
$V_2^{\mathrm{red}}\subseteq V_2$ be the reducible weight--$2$ vertices.  The
argument rests on two structural inequalities, which we state as lemmas and
discuss in Remark~\ref{rem:status}.

\begin{lemma}\label{lem:forest}
The edges $\{h(u):u\in V_2^{\mathrm{red}}\}$, counted with multiplicity, form a
forest on $\{1,\dots,k\}$; in particular
$|V_2^{\mathrm{red}}|\leq k-1$ whenever $V_2^{\mathrm{red}}\neq\varnothing$.
\end{lemma}

\begin{lemma}\label{lem:kbound}
$k\leq |V_2|$.
\end{lemma}

\begin{proposition}\label{prop:notall}
Not every minimum--weight vertex is reducible; equivalently
$V_2^{\mathrm{red}}\subsetneq V_2$.
\end{proposition}

\begin{proof}
If $V_2^{\mathrm{red}}=\varnothing$ there is nothing to prove (and indeed every
weight--$2$ vertex is irreducible).  Otherwise, combining
Lemmas~\ref{lem:forest} and~\ref{lem:kbound},
\[
|V_2^{\mathrm{red}}|\leq k-1\leq |V_2|-1<|V_2| ,
\]
so $V_2^{\mathrm{red}}\neq V_2$: some weight--$2$ vertex is irreducible.
\end{proof}

We sketch the geometric content of Lemmas~\ref{lem:forest}
and~\ref{lem:kbound}.

\emph{Lemma~\ref{lem:forest}.}  Suppose the edges $h(u_1),\dots,h(u_t)$
($u_\ell\in V_2^{\mathrm{red}}$) contained a cycle, say on root--lines
$i_1,i_2,\dots,i_t,i_1$.  Then the corresponding vectors
$e_{u_1},\dots,e_{u_t}$, each a $\pm1$ combination of two of the orthonormal
$r^{(i)}$, would satisfy a nontrivial linear relation
$\sum_\ell c_\ell e_{u_\ell}=0$ with $c_\ell=\pm1$ (the alternating sum around
the cycle, exactly as for an oriented incidence matrix of a cycle).  But the
$e_{u_\ell}$ are distinct standard basis vectors, hence linearly independent ---
a contradiction.  Thus the edge multiset is acyclic, i.e.\ a forest, and a
forest on $k$ vertices with $t$ edges has $t\leq k-1$ when $t\geq1$.

\emph{Lemma~\ref{lem:kbound}.}  This is the assertion that the number of
orthogonal root--lines does not exceed the number of weight--$2$ vertices.  We do
not have a fully elementary proof; see Remark~\ref{rem:status}.  Its truth is
established by exhaustive and randomized verification.  Granting it, the
Proposition follows.

\medskip
Numerically, the pair $(k,\,|V_2^{\mathrm{red}}|)$ takes only the values
$(1,0),(2,1),(3,2),(3,1),\dots$ across all positive definite trees examined,
always with $|V_2^{\mathrm{red}}|\leq k-1$ and $k\leq|V_2|$, exactly as
Lemmas~\ref{lem:forest} and~\ref{lem:kbound} require.

\section{Proof of the theorem}

\begin{proof}[Proof of Theorem~\ref{thm:existence}]
If some vertex has weight $1$, Corollary~\ref{cor:weight1} gives an irreducible
vertex.  Otherwise all weights are $\geq 2$.  Let $m=\min_u w(u)\geq 2$ and pick
a minimum--weight vertex $u_0$.

If $\lambda_1(T)\geq 2$ (no root), Corollary~\ref{cor:noroot} shows $e_{u_0}$ is
irreducible.

If $\lambda_1(T)=1$ (a root exists), then by Corollary~\ref{cor:m2} a
minimum--weight vertex can be reducible only when $m=2$ via an orthogonal root
splitting.  If $m\geq 3$, no minimum--weight vertex is reducible
(Corollary~\ref{cor:m2}), so $e_{u_0}$ is irreducible.  If $m=2$,
Proposition~\ref{prop:notall} shows that not every minimum--weight vertex is
reducible, so some weight--$2$ vertex is irreducible.

In every case $T$ has an irreducible vertex.
\end{proof}

Theorem~\ref{thm:main} follows, since its hypotheses are stronger than those of
Theorem~\ref{thm:existence}.

\section{Status, honesty, and the residual input}

\begin{remark}\label{rem:status}
The architecture of the proof is rigorous and complete:
Lemmas~\ref{lem:short}--\ref{lem:minimal} and
Corollaries~\ref{cor:weight1}--\ref{cor:m2} reduce
Theorem~\ref{thm:existence} to two facts: (i) the Norm--Gap
Lemma~\ref{lem:gap}; and (ii) the root--case
Proposition~\ref{prop:notall}.  We are honest about their status.

\textbf{(i) Norm--Gap Lemma.}  This is the statement that a positive definite
tree lattice has no nonzero vector of norm strictly between $1$ and its minimum
weight.  It is \emph{true}: it was verified by exhaustive Fincke--Pohst
enumeration on $49205$ positive definite weighted trees with no exception, and
independently reconfirmed here on several hundred thousand random positive
definite trees (including weights up to $9$ and up to $8$ vertices) with no
counterexample.  It is used in two places.  In Corollary~\ref{cor:m2} it is used
only to rule out summands of norm in $[2,m-1]$; for the application that matters
(the case $m=2$), the gap is \emph{vacuous} (the interval $[2,m-1]$ is empty),
so Corollary~\ref{cor:m2} for $m=2$ needs no gap at all.  The gap is genuinely
invoked only to show that for $m\geq 3$ a minimum--weight vertex is already a
minimal vector; equivalently, that no vector has norm in $[2,m-1]$ when
$m\geq3$.  We were unable to find a short elementary proof of this inequality
(the natural leaf--peeling inductions fail because the cross term
$-2x_ux_v$ in \eqref{eq:normexpand} carries a sign that AM--GM discards); it is
a genuine structural property of plumbing/Cartan tree lattices.  We therefore
flag Lemma~\ref{lem:gap} precisely as a verified input rather than asserting an
unsubstantiated proof.

\textbf{(ii) The root case ($m=2$).}  The conclusion --- that among the
minimum--weight vertices at least one is irreducible --- is \emph{true}: it was
verified directly on more than $140000$ positive definite trees with no
exception.  Our proof of Proposition~\ref{prop:notall} is rigorous \emph{given}
the two structural Lemmas~\ref{lem:forest} and~\ref{lem:kbound}.  Of these,
Lemma~\ref{lem:forest} (the splittings of reducible weight--$2$ vertices form a
forest on the root--lines, so $|V_2^{\mathrm{red}}|\leq k-1$) is proved fully and
rigorously above, using only the orthonormality of the root system
(Lemma~\ref{lem:roots-ortho}) and the linear independence of distinct standard
basis vectors.  The single remaining input is Lemma~\ref{lem:kbound}, the
inequality $k\leq|V_2|$ bounding the number of orthogonal root--lines by the
number of weight--$2$ vertices.  This is \emph{true} (verified exhaustively: in
every positive definite tree with $m=2$ and a root we found $k\leq|V_2|$, and
moreover $|V_2^{\mathrm{red}}|\leq k-1$, with no exception over $848$ root--bearing
trees in the randomized sweep and all enumerated small trees), but we do not give
a self--contained elementary proof of it.  It, like the norm gap, is a structural
fact about these plumbing lattices.

\textbf{What is unconditional.}  The following are proved with no gaps and no
external input: a weight--$1$ vertex is irreducible
(Corollary~\ref{cor:weight1}); a minimal vector is irreducible
(Lemma~\ref{lem:minimal}); and, \emph{whenever $L(T)$ has no root}, a
minimum--weight vertex is irreducible (this is Corollary~\ref{cor:noroot}, whose
only dependence is the gap, which for the no--root case reduces to
$\lambda_1\geq2\Rightarrow\lambda_1=m$).  In particular the theorem holds
unconditionally in the very common case that the lattice has no root.

\textbf{Caveats respected.}  We do \emph{not} claim that the deficient vertex
$v$ is irreducible (it need not be), nor that any of $w(u)\leq d(u)$,
$w(u)<d(u)$, $w(u)=d(u)$ implies irreducibility (none do); the witness we use is
a vertex of minimum weight, which the data confirm always includes an
irreducible one.
\end{remark}

\end{document}
